\section{Inference and Finite-Sample Audit}\label{app:inference}

\subsection{Stochastic targets}

The principal wild-score interval treats occupation as the shock cluster and conditions on the constructed stocks, exposure labels, occupational bridge, and retained support.  It accommodates arbitrary serial dependence within an occupation but not unrestricted dependence across occupations.  A 22-family Webb-weight exercise instead allows common shocks within broad SOC families.  A separate household/sample-unit multiplier rebuilds cell stocks from microdata, preserves all observed records of the sampled unit across months, and assigns the same multiplier to every fractional descendant of a routed source record.  That exercise describes repeated-sample sensitivity conditional on released CPS final weights; the public Basic Monthly extracts do not supply the design variables needed to call it complete CPS design-based inference.

These variances are not mechanically added.  Occupation-shock and household-resampling procedures can contain overlapping components, so summing them without a derived orthogonal decomposition could double count uncertainty.  They are reported as separate sensitivity targets.

\subsection{Score and bootstrap implementation}

The conditional likelihood retains one-sided zero cells and drops only occupation--month cells with zero total stock.  Scores use the fitted grouped-binomial objective, and nuisance-adjusted target scores use the Schur-complement projection implied by the fitted information matrix.  The same 9,999 Rademacher occupation multipliers are reused for paired comparisons.  Each draw first forms the two refitted or score-updated targets and then their difference; intervals for movement therefore retain covariance.  The finite-cluster correction, random seed, multiplier matrix hash, convergence checks, and score-conservation checks are written to the execution receipts.

For the reconstructed pooled comparison, the analytic standard error is 0.04517 and the 9,999-draw 95-percent interval is $[-0.2206,-0.0437]$.  After SOC2-by-calendar-month conditioning, the coefficient is $-0.0217$, with standard error 0.07132 and interval $[-0.1607,0.1173]$.  The paired movement is 0.11043, with standard error 0.05189 and interval $[0.0107,0.2102]$.  The corresponding two-sided normal-theory paired $\mathrm{MDE}_{80}$ is 0.14539.  The detected movement remains smaller than this planning quantity; that is not a contradiction because an MDE is a power summary, not an acceptance threshold for an observed draw.

\subsection{Broad-family and household sensitivity}

The few-family exercise uses 22 broad-family weights and reports baseline, conditioned, and paired targets from common draws.  Its role is sensitivity to cross-occupation dependence, not a declaration that SOC2 is the uniquely correct cluster.  The household/sample-unit exercise performs 199 positive-weight full refits.  Under reconstructed exposure labels, its percentile interval is $[-0.1794,-0.0800]$ for the pooled coefficient, $[-0.1119,0.0560]$ for the family-conditioned coefficient, and $[0.0370,0.1722]$ for their movement.  Labels and support remain fixed because the target is repeated-sample uncertainty conditional on the preperiod construction; regenerating the generated treatment would define a broader and much more expensive target.

The source does not expose official Basic Monthly replicate weights or all confidential design identifiers.  Household identifiers are therefore a dependence-preserving public-data approximation, not an ASEC replication-weight procedure.  Link-specific flow analyses additionally report person- and household-cluster sensitivities using their own risk sets.

\subsection{Calendar-time HAC and stress validation}

The time-HAC audit operates on nuisance-adjusted score sums and implements occupation clustering plus an aggregate elapsed-calendar HAC term minus within-occupation elapsed-calendar HAC terms.  Calendar lags are measured in months, so the missing October 2025 observation does not make September and November adjacent.  Under rebuilt exposure labels, lag-16 standard errors are 0.04449 for the pooled target, 0.05522 for the family-conditioned target, and 0.04501 for their paired movement.  None of the 15 rebuilt five-parameter covariance matrices across the three targets and five lags is indefinite, and no projection or eigenvalue clipping is applied.  Indefinite positive-lag matrices from the superseded historical-treatment/SOC2-post audit are retained only as labeled implementation history.

The completed finite-sample audit uses 11 declared designs, 799--1,599 retained outer replications per design, 9,999 common multiplier draws per joint refit, and independently computed pseudo-targets.  Absolute coefficient bias is at most 0.00102 log point in the empirical observed-effect design.  Interval performance is target-specific: occupation-Rademacher coverage is 0.865 for the pooled projection, 0.921 for the family-month projection, and 0.889 for their paired movement.  Family-Rademacher coverage is respectively 0.990, 0.895, and 0.992.  Rademacher and Webb weights give nearly identical coverage at a fixed cluster level, so changing the multiplier law does not repair the cluster-level mismatch.

The earlier 26.7-percent and 11.3-percent rejection figures are withdrawn rather than reinterpreted.  In the clean family-variance-zero design, the family-month pseudo-target is zero to numerical precision; its five-percent rejection rates are 6.1 percent under occupation Rademacher weights, 8.9 percent under family Rademacher weights, and 5.6 percent for the cross-fit full-refit oracle.  In the empirically calibrated structural-null design, removing the structural exposure coefficient still leaves a pooled pseudo-target of $-0.11035$ and a paired pseudo-target of 0.11035 because the data-generating process retains the observed fitted family-by-month surface.  The near-nominal oracle then rejects zero for the pooled projection 96.5 percent of the time.  This is a restatement of the broad-family decomposition built into the calibration, not independent evidence that composition caused the empirical gap.  The simulations therefore limit both the inferential and structural interpretation: none of the feasible cluster procedures is adequate for every target and significance of the pooled projection does not by itself test a structural exposure effect.
